\bprob[title={DoCarmo, Riemannian Geometry, page 83 problem 7}]

    Let $M$ be a Riemannian manifold, and $p\in M$.
    Show that there is a neighborhood $U$ of $p$ in $M$ and vector fields $E_1,\dots,E_n\in\X(U)$ forming an orthonormal local frame such that
    $\nabla_{E_i}E_j|_p=0$ for all $i,j$.
    Such a frame is called a {\emphcolor local geodesic frame}.

\eprob

Let $B$ be a geodesic ball around $p$, so of the form $B=\exp_p(B_\epsilon(0))$ for an open ball $B_\epsilon(0)\subseteq T_pM$ contained in the domain of $\exp_p$.
Furthermore let $v_1,\dots,v_n\in T_pM$ form an orthonormal basis.
Now for every $q\in B$ consider the radial geodesic from $p$ to $q$, so if $q=\exp_p(w)$ for $w\in B_\epsilon(0)$ the radial geodesic is given by $\gamma_q\colon t\mapsto\exp_p(tw)$.
Then define $E_i|_q$ to be the parallel transport of $v_i\in T_pM$ along $\gamma_q$:
$$ E_i|_q = P^{\gamma_q}_{0,1}(v_i) . $$

Parallel transport along a curve is a linear isometry.
If $\gamma$ is a curve, $v,u\in T_pM$, and $X,Y$ are their parallel transports along $\gamma$, then by definition $P^\gamma_{0,1}(v)=X(1)$ so
$$ \iprod{P^\gamma_{0,1}v,P^\gamma_{0,1}(u)} = \iprod{X(1),Y(1)} . $$
On the other hand $\iprod{X(t),Y(t)}$ is constant along $\gamma$, since
$$ \drv t\iprod{X(t),Y(t)} = \iprod{D_tX,Y} + \iprod{X,D_tY} = 0 . $$
And so
$$ \iprod{P^\gamma_{0,1}v,P^\gamma_{0,1}u} = \iprod{X(1),Y(1)} = \iprod{X(0),Y(0)} = \iprod{v,u} , $$
as desired.
Since $E_i|_q$ is the parallel transport of $v_i$ along $\gamma_q$, this means the $E_i|_q$ remain orthonormal.

So we just need to show that the $E_i$ are smooth, and $\nabla_{E_i}E_j|_p=0$.
Consider the definition of $E_i$: for $q\in B$ we compute $w=\exp_p^{-1}(q)$ and then the parallel transport of $v_i$ along $\gamma=\exp_p(tw)$.
We will show that the parallel transport of $v=v_i$ along $\gamma$ is smooth in $w$.
The parallel transport is given by $X(1)$ where $X$ solves the ODE
$$ \pdvof{X^i}t(t) + X^i(t)\dot\gamma^j(t)\Gamma^k_{ij}(\gamma(t)) = 0 , $$
since $\gamma$ depends smoothly on $w$, the solution to this ODE depends smoothly on $w$.
And $w$ depends smoothly on $q$ since $\exp_p^{-1}$ is a diffeomorphism, so in total $X$ depends smoothly on $q$.
And thus $E_i|_q=X(1)$ depends smoothly on $q$, so $E_i$ is smooth.

To show that $\nabla_{E_i}E_j|_p=0$, we have that
$$ \nabla_{E_i}E_j|_p = \nabla_{E_i(p)}E_j = \nabla_{v_i}E_j . $$
Now let $\alpha_i$ be the geodesic starting at $p$ with initial velocity $v_i$, so $\alpha_i=\exp_p(tv_i)$.
So this is equal to
$$ = \nabla_{\alpha'_i(0)}E_j = D_t|_{t=0}(E_j\circ\alpha_i) , $$
with the covariant derivative taken along $\alpha_i$.
Now, $E_j(\alpha_i(t))$ is the parallel transport of $v_j$ along $\gamma_{\alpha_i(t)}=\alpha_i$, and thus it is by definition parallel along $\alpha_i$ and its covariant
derivative is zero, as desired.

\bprob

    \benum
        \item Let $O(1,n)$ denote the subgroup of linear transformations of $\bR^{n+1}$ which preserve the Lorentz bilinear form.
        Let $O_+(1,n)$ denote the set of matrices $A^i_j$ in $O(1,n)$ with $A^0_0>0$.
        Show that $O_+(1,n)$ is a subgroup of $O(1,n)$ and its transformations restrict to isometries of $H^n_k$.

        \item Determine the geodesics in the hyperboloid model of hyperbolic space as follows.
        Let $P$ be a two-dimensional linear subspace of $\bR^{n+1}$ containing the point $p\in H^n_k$.
        Show that there is an element of $O_+(1,n)$ which fixes $P$ and acts non-trivially on $\bR^{n+1}-P$.
        Conclude that all geodesics through $p$ can be obtained as intersections of $P\cap H^n_k$ for some $P$ as above.

        \item Show that the geodesics of point (2) are sent by the isometry $h\colon H^n_1\to D^n$ (where $D^n$ is the Klein model of hyperbolic space) to straight lines.

        \item Show that geodesics in the Poincar\'e model of hyperbolic space are circles perpendicular to the boundary of the disk.

        \item Show by example that angles measured with respect to the Klein metric on the disk are not the same as with respect to the Euclidean metric on the disk.
        On the other hand, show that angles on the Poincar\'e disk are the same as the Euclidean metric.
    \eenum

\eprob

\benum
    \item Let us denote by $Q=\diag(-1,1,\dots,1)$ the bilinear form, so $(x,y)=x^\top Qy$
    If we write $x=(a,u)$ and $y=(b,v)$ then we see $(x,y)=-ab+\iprod{u,v}$ where $\iprod{\bul,\bul}$ denotes the Euclidean product on $\bR^n$.

    Then $A\in O(1,n)$ iff
    $$ (Ax,Ay) = x^\top A^\top QAy = x^\top Qy = (x,y) , $$
    i.e.\ iff $A^\top QA=Q$.
    Since $Q$ is its own inverse, clearly $O(1,n)$ consists of invertible matrices.
    Clearly then $O(1,n)$ is closed under multiplication and contains the identity, and by multiplying $A^\top QA=Q$ by $(A^{-1})^\top$ and $A^{-1}$ we see it is closed under inversion, so it is a group.
    Furthermore we see that $A\in O(1,n)$ iff
    $$ A^{-1} = QA^\top Q . $$
    Then by taking the transpose of both sides we see that this implies
    $$ (A^\top)^{-1} = (A^{-1})^\top = QAQ , $$
    so that $O(1,n)$ is further closed under transposition.

    Now to show that $O_+(1,n)$ is a subgroup, we first note that it clearly contains the identity.
    And if $A\in O_+(1,n)$ then since $A^{-1}=QA^\top Q$, we see that $(A^{-1})^0_0=A^0_0>0$ so $O_+(1,n)$ is closed under inversion.
    And suppose that $A\in O(1,n)$ and let
    $$ A = \pmatrix{a & u^\top\cr v & *} . $$
    Then, letting $e_i$ be the standard basis vectors,
    $$ -1 = (e_0,e_0) = (Ae_0,Ae_0) = ((a,v),(a,v)) = -a^2 + \norm v^2 \implies a^2 = \norm v^2 + 1 , $$
    where the norm is the Euclidean norm.
    In particular we have that $a>\norm v$.
    Since $O(1,n)$ is closed under transposition, we see that $a^2=\norm u^2+1$ as well.
    If $A\in O_+(1,n)$ so that $a>0$ we see that $a>\norm u,\norm v$.

    Now suppose we have
    $$ A = \pmatrix{a & u^\top\cr v & *}, B = \pmatrix{b & z^\top\cr w & *} \in O_+(1,n) , $$
    we want to show that $AB\in O_+(1,n)$.
    We have
    $$ AB = \pmatrix{ab + u^\top w & *\cr * & *}, $$
    so we want to show that $ab+u^\top w>0$.
    By Cauchy-Schwarz we have
    $$ ab + u^\top w \geq ab - \norm u\norm w > ab - ab  = 0, $$
    where the final inequality is because we showed $a>\norm u$ and $b>\norm w$ above.
    Thus $O_+(1,n)$ is indeed a subgroup.

    By definition $O(1,n)$, and thus $O_+(1,n)$, are isometries of Lorentzian space.
    Since $H^n_k$ is an embedded submanifold of Lorentzian space, if we can show that maps of $O_+(1,n)$ restrict to it then we will have finished.
    We want to show that if $A\in O_+(1,n)$ and $x\in H^n_k$ then $Ax\in H^n_k$, so $(Ax)^0>0$ and $Q(Ax)=1/k$.
    Since $A$ is an isometry, $Q(Ax) = Q(x) = 1/k$ so we need only show that $(Ax)^0>0$.

    Since $x\in H^n_k$ we have that $x^0>0$ and $Q(x)<0$.
    Let $x=(x_0,y)$ so $Q(x)=-x_0^2+\norm y^2<0$ meaning $x_0>\norm y$.
    Now let
    $$ A = \pmatrix{a & u^\top\cr v & *}, $$
    then $(Ax)^0=ax_0+u^\top y$ and by Cauchy-Schwarz,
    $$ (Ax)^0 = ax_0 + u^\top y \geq ax_0 - \norm u\norm y > ax_0 - ax_0 = 0 , $$
    as desired.

    \item First we claim that $A\in O(1,n)$ is in $O_+(1,n)$ iff it preserves the sign of the first compomnent of any vector $x$ with $(x)<0$.
    If $A\in O_+(1,n)$ then in the previous point we showed that if $x^0>0$ and $(x)<0$ then $(Ax)^0>0$ as well.
    And if $x^0<0$ and $(x)<0$ then $-x^0>0$ and $(-x)=(x)<0$ so $(-Ax)^0>0$ so $(Ax)^0<0$, so $A$ preserves the sign of the first component all such vectors.
    And conversely if $A$ preserves the sign of all such vectors, then in particular it does for $e_0$, and $(Ae_0)^0=A^0_0>0$ so $A\in O_+(1,n)$.

    We also claim that $A\in O_+(1,n)$ iff it preserves the sign of the first component of {\it a single\/} $x$ with $(x)<0$.
    This is because if $A\notin O_+(1,n)$ then $-A\in O_+(1,n)$ (since $(Ae_0)=-1$ we must have $A^0_0\neq0$), so $-A$ must preserve the sign of $x$.
    But $(-Ax)^0=-Ax^0$ which differs from the sign of $x$ since $A$ preserves it, a contradiction.

    Let $P\leq\bR^{n+1}$ be two-dimensional containing $p$.
    Letting $x_0$ be the normalization of $p$ we can extend $x_0$ to an orthonormal basis $x_0,x_1$ of $P$.
    This in turn we extend to an orthonormal basis $x_0,\dots,x_n$ of $\bR^{n+1}$.
    By Sylvester's law of inertia, $(x_i,x_j)=\pm\delta_{ij}$ with the negative sign occuring only when $i=j=0$.
    Now restricting the bilinear form to $P^\perp=\lspan\set{x_2,\dots,x_n}$ makes it an inner product since $(x_i,x_j)=\delta_{ij}$ for $i,j\geq2$.
    So define $A$ on $P^\perp$ to be a unitary transformation with no fixed points, and to be the identity on $P$.
    This preserves the bilinear form, so it is in $O(1,n)$.
    Furthermore it preserves the sign of $x_0$, and since $(x_0)<0$, by above this means $A\in O_+(1,n)$.

    Thus by definition $A\in O_+(1,n)$ fixes $P$ and for every $x\in\bR^{n+1}-P$ we can write $x=x'+x''$ for $x'\in P,x''\in P^\perp$,
    then $Ax=x'+Ax''\neq x'+x''$, so every point outside of $P$ is acted on non-trivially.

    Now if $\gamma$ is a geodesic in $H^n_k$ passing through $p$ with velocity $v$ at $p$.
    So $\gamma(t_0)=p$ and $\gamma'(t_0)=v\in T_pH^n_k$, which means that $(v,p)=0$.
    Let $P$ be spanned by $p,v$ (since $(v,p)=0$ they are linearly independent so $P$ is two-dimensional).
    Let $A\in O_+(1,n)$ be as above, so $Ap=p,Av=v$.
    As $A$ is an isometry, $A\circ\gamma$ is a geodesic of $H^n_k$, and
    $$ A\circ\gamma(t_0) = p,\qquad (A\circ\gamma)'(t_0) = dA_p\circ\gamma'(t_0) = v,  $$
    since as a matrix, $A$'s differential can be identified with $A$.
    So by uniqueness, $A\circ\gamma=\gamma$, and since $A$'s fixed points are $P$ this means that $\gamma$'s image lies in $P\cap H^n_k$.

    Now, $P\cap H^n_k$ are vectors $\alpha p+\beta v$ such that
    $$ (\alpha p+\beta v) = (\alpha p) + (\beta v) = -r^2\alpha^2 + \beta^2(v) = -r^2 , $$
    where $r^2=-1/k$.
    These are points along a hyperbola, and thus $\gamma$'s image is the image of an open interval in this hyperbola.
    Its image must be the entire hyperbola, since otherwise by parameterizing the hyperbola ($\alpha(t)=\cosh(t)$ and $\beta(t)=r\sinh(t)/\sqrt{(v)}$ since $(v)>0$) we can take the
    first $t$ not in $\gamma$'s image.
    Letting $q=\alpha(t)p+\beta(t)v\in H^n_k\cap P$ not be in $\gamma$'s image, and suppose $\gamma$ is defined on $(a,b)$.
    But then $\gamma(t)\to q$ as $t$ approaches the endpoint of $\gamma$'s domain, wlog $b$.

    Let $W$ be a uniformly normal neighborhood of $q$, then eventually $\gamma(t)$ is in $W$.
    So eventually $\gamma(t)$ is contained in $W$, let $t_0$ be such that $\gamma(t)\in W$ for $t\geq t_0$, and $s=\gamma(t_0)$.
    Then $\gamma(t)=\exp_s(tw)$ for some $w\in T_sH^n_k$.
    Since $\gamma(t)\to q$ this means that $\exp_s(tw)\to q$, and $q\in W$ so $q=\exp_s(u)$ for some tangent vector $u$ (since $W$ is uniformly normal).
    Thus $tw\to u$, and thus $u=bw$, so we can extend $\gamma$ to its value at its endpoint to be $\gamma(b)=q$ and this remains smooth (since it is given by $\exp_s(tw)$).
    This is a contradiction to $\gamma$'s maximality.
    Thus $\gamma$'s image is the entirety of $P\cap H^n_k$.

    \item We want to show that $h$ maps $P\cap H^n_{-1}$ to a straight line.
    Elements of this intersection are vectors $(p_0,p)\in H^n_{-1}$ which are also in $P$.
    Since a $k$-dimensional subspace of $n$-dimensional space can be written as the kernel of an $(n-k)\times n$ matrix, and thus can be written as solutions to a system of
    $n-k$ homogeneous linear equations, elements of $P$ are solutions to a system of $n-1$ linear equations.
    Now $(p_0,p)$ solves a linear equation $\alpha p_0+\beta^ip_i=0$ iff $h(p_0,p)=p/p_0=q$ solves $\beta^iq_i=-\alpha$.
    These equations define an affine subspace of $\bR^n$, thus $h(P\cap H^n_{-1})=h(P)\cap D^n$ is the intersection of $D^n$ with an affine subspace of $\bR^n$.
    Since $h(P\cap H^n_{-1})$ is the image of a geodesic and thus a curve, it is one-dimensional and thus must be the intersection of $D^n$ with a line, as desired.

    \item First we consider the two-dimensional case, and we consider the stereographic projection $f\colon H^n_{-1}\to D^n$ given by
    $$ f(p_0,p) = \frac p{1+p_0} . $$
    This is an isometry, and thus maps geodesics to geodesics.
    Since geodesics in $H^n_{-1}$ are the solutions to $n-1$ linear equations, we first restrict to the two-dimensional case, so geodesics are solutions
    to a single linear equation $-\alpha p_0+\beta^ip_i=0$, rewriting it as $\alpha p_0=\iprod{\beta,p}$ where the inner product is Euclidean.
    So given a geodesic parameterized by $p$, $u=f(p)$ is a geodesic.

    If $\alpha=0$ then $\iprod{\beta,p}=0$ iff $\iprod{\beta,f(p_0,p)}=0$, so $f(p_0,p)$ forms a straight line orthogonal to $\beta$.
    This is a straight line through the origin and thus orthogonal to the boundary of the disk.

    If $\alpha\neq0$ then by normalizing $\alpha$ and $\beta$ we can assume $\alpha=1$ so our equation is $p_0=\iprod{\beta,p}$.
    If we consider $f^{-1}$ given by
    $$ f^{-1}(u) = \parens{\frac2{1-\norm u^2}-1,\ \frac{2u}{1-\norm u^2}} , $$
    then points on the geodesic are the points $f^{-1}(u)$ where
    $$ \frac2{1-\norm u^2}-1 = \iprod{\beta,\frac{2u}{1-\norm u^2}} \iff \norm u^2 - 2\iprod{\beta,u} + 1 = 0 . $$
    This can be written as
    $$ \norm{u-\beta}^2 = \norm\beta^2 - 1 . $$
    If $\norm\beta^2\leq1$ then this has either no solution or the solution is $u=\beta$ when $\norm\beta=1$ in which case this lies on the boundary of the disk and not within the disk
    itself.

    Otherwise this defines a circle of radius $\sqrt{\norm\beta^2-1}$ with center $\beta$, so the geodesic is the intersection of that circle with $D^n$.
    Now consider the intersection of this circle with the boundary of $D^n$, so the $u_0$s where $\norm{u_0-\beta}^2=\norm\beta^2-1$, $\norm{u_0}=1$.
    To show that the geodesic is orthogonal to the disk, we show that $0,u_0,\beta$ form a right triangle, that is we want to show $\iprod{u_0,u_0-\beta}=0$.
    We have that 
    $$ \norm\beta^2 - 1 = \norm{u_0}^2 - 2\iprod{u_0,\beta} + \norm\beta^2 \implies \iprod{u_0,\beta} = 1 , $$
    and thus
    $$ \iprod{u_0,u_0-\beta} = \norm{u_0}^2 - \iprod{u_0,\beta} = 1 - \iprod{u_0,\beta} = 0 , $$
    as desired.

    In the higher-dimension case, a geodesic of $H^n_{-1}$ is determined by a hyperbola on a two-dimensional affine subspace, say $\alpha p+\beta q=v$.
    Writing our coordinates as $(p_0,p_1,\dots,p_n)$, we can perform an orthogonal transformation on $p_1,\dots,p_n$ (which does not change the metric clearly) to map this
    affine subspace down to one existing in the coordinates $(p_0,p_1,p_2,0,\dots,0)$.
    Then we are in the same situation as the two-dimensional case.

    \item To measure the angle between two geodesics $\gamma_1,\gamma_2$ intersecting at a point $p=\gamma_1(0)=\gamma_2(0)$, we compute the angle between
    $\gamma'_1(0),\gamma'_2(0)\in T_pM$ using the Riemannian metric.
    In the Klein model, to geodesic connecting any two points is the same as in Euclidean space.
    Thus to measure the angle between any three points, we are measuring the angle between the same tangent vectors.

    If we take an arbitrary point, $p\in D^n$ and two tangent vectors $v,u\in T_pD^n$ then the angle between them is given by
    $$ \cos\theta = \frac{\iprod{u,v}}{\norm u\norm v} . $$
    Let us take $p=(1/2,1/2)\in D^2$ and $u=e_1,v=e_2$ be the standard basis vectors.
    Then in the Euclidean metric, $u$ and $v$ are orthogonal and thus $\cos\theta=0$.
    But in the Klein metric,
    $$ k_{12} = \frac{p_1p_2}{(1-\norm p^2)^2} \neq 0 , $$
    and so $e_1,e_2$ are not orthogonal, so the angle between them is not $90^\circ$, different than in the Euclidean case.
\eenum

\bprob

    Let $M\subseteq\bR^{n+1}$ be the graph of the smooth function $f\colon\bR^n\to\bR$,
    $$ M = \set{(x_0,\bar x)}[x_0=f(\bar x)] . $$
    Give $\bR^{n+1}$ either the Euclidean or Lorentzian metric, $G_{ij}=\pm\delta_{ij}$ where the negative sign occurs only potentially when $i=j=0$.
    Let $g$ be the pullback of $G$ to $M$, and let $\x\colon\bR^n\to M$ be the coordinates
    $$ \x(x_1,\dots,x_n) = (f(x_1,\dots,x_n),x_1,\dots,x_n) . $$
    Write $f_i=\ipdv{x^i}f$.
    \benum
        \item Show that, with respect to $\x$, $g_{ij}=\delta_{ij}\pm f_if_j$.
        Conclude that in a small enough neighborhood $U$ of a critical point of $x$, $g$ defines a Riemannian metric on $\x(U)$.
        We implicitly make this restriction for the rest of the problem.

        \item With respect to $\x$ show that $g$'s Christoffel symbols are
        $$ \Gamma^k_{ij} = \pm f_{ij}f_\ell g^{\ell k} . $$

        \item Show that for a critical point $x$ of $f$, we have
        $$ R_{ijk\ell} = \pm(f_{ik}f_{j\ell} - f_{i\ell}f_{jk}) . $$
    \eenum

\eprob

\benum
    \item Let us denote by $\partial_i$ the coordinate frame of Euclidean space, then $M$'s coordinate frame with respect to $\x$ is given by $\tilde\partial_i=d\x(\partial_i)$.
    As $M\subseteq\bR^{n+1}$ is embedded, we can compute this relative to the coordinate frame of $\bR^{n+1}$ (also denoted $\partial_i$),
    $$ \tilde\partial_i = \pdvof{\x^j}{\partial_i}\partial_j = f_i\partial_0 + \partial_i . $$
    Then since $g$ is the pullback of $G$, we can compute the inner product in Euclidean space, so we have
    $$ g_{ij} = \iprod{\tilde\partial_i,\tilde\partial_j} = \iprod{f_i\partial_0 + \partial_i, f_j\partial_0 + \partial_j} = f_if_jG_{00} + G_{ij} = \delta_{ij} \pm f_if_j , $$
    as desired.
    A critical point $x$ of $f$ is one where $f_i(x)=0$ for all $i$, and by smoothness there is a neighborhood of $x$ where $\abs{f_if_j}<1$ for all $i,j$.
    In such a neighborhood $g_{ij}>0$ and thus $g$ is Riemannian, as the pullback of $G$ it is already symmetric.

    \item We recall
    $$ \Gamma^k_{ij} = \frac12g^{k\ell}(\tilde\partial_ig_{j\ell} + \tilde\partial_jg_{i\ell} - \tilde\partial_\ell g_{ij}) . $$
    And $\tilde\partial_i=d\x(\partial_i)$ acts on smooth functions by
    $$ \tilde\partial_if = \partial_i(f\circ\x) , $$
    thus
    $$ \tilde\partial_ig_{j\ell} = \partial_i(\delta_{ij} \pm f_jf_\ell) = \pm f_{ij}f_\ell \pm f_jf_{i\ell} . $$
    Plugging this in we see
    $$ \Gamma^k_{ij} = \pm\frac12g^{k\ell}(f_{ij}f_\ell + f_jf_{i\ell} + f_{ij}f_\ell + f_if_{j\ell} - f_{i\ell}f_j - f_if_{j\ell}) = \pm f_{ij}f_\ell g^{k\ell} , $$
    as desired ($g^{k\ell}$ is symmetric).

    \item If $x$ is a critical point of $f$, then $f_i=0$ for all $i$.
    Thus $\Gamma^k_{ij}=0$ and $g_{ij}=\delta_{ij}$ for all $i,j,k$.
    We know the components of Riemannian curvature are given by
    $$ R_{ijk\ell} = -g_{\ell m}(\tilde\partial_i\Gamma^m_{jk} - \tilde\partial_j\Gamma^m_{ik} + \Gamma^p_{jk}\Gamma^m_{ip} - \Gamma^p_{ik}\Gamma^m_{jp}) = \tilde\partial_j\Gamma^\ell_{ik} - \tilde\partial_i\Gamma^\ell_{jk} . $$
    And
    $$ \tilde\partial_i\Gamma^\ell_{jk} = \pm\partial_i(f_{jk}f_mg^{m\ell}) = \pm(f_{ijk}f_mg^{m\ell} + f_{jk}f_{im}g^{m\ell} + f_{jk}f_m\partial_ig^{m\ell}) = \pm f_{jk}f_{im}\delta_{m\ell} = \pm f_{jk}f_{i\ell} . $$
    Thus plugging in we see
    $$ R_{ijk\ell} = \pm(f_{ik}f_{j\ell} - f_{jk}f_{i\ell}) . $$
\eenum

\def\hol{{\rm Hol}}

\bprob

    Let $M$ be a Riemannian manifold, and $x,y\in T_pM$ tangent vectors.
    Let $\gamma_t$ be the loop in $M$ based at $p$ obtained by applying $\gamma_t$ to the parallelogram in $T_pM$ spanned by $tX,tY$.
    Parallel transport along a closed loop is called {\emphcolor holonomy}, and thus this defines an operator
    $$ \hol_{\gamma_t}\colon T_pM\to T_pM . $$
    Show that
    $$ R(X,Y) = \evdrv t_{t=0}\hol_{\gamma_{\sqrt t}} . $$

\eprob

Let us define the smooth family of curves $\Gamma\colon I\times I\to M$ by $\Gamma(s,t)=\exp(sx+ty)$ for $I$ an open interval containing $0$ small enough so that this is defined.
We have $\partial_s\Gamma(0,0)=x$ and $\partial_t\Gamma(0,0)=y$.
The holonomy operator of $\gamma_a$ is obtained via parallel transport along $\Gamma(s,0)$ then $\Gamma(a,t)$ then $\Gamma(a-s,a)$ then $\Gamma(0,a-t)$.

Now in general if we have a curve $\gamma\colon[0,a]\to M$, then we can define its reverse $-\gamma\colon[0,a]\to M$ by $t\mapsto\gamma(a-t)$.
Then if $X$ is parallel along $\gamma$ then $X^-$ is parallel along $-\gamma$ where $X^-(t)=X(a-t)$.
This is because for $X\circ\gamma$ extendible $(X\circ\gamma)^-=X\circ(-\gamma)$ and then
$$ D_t^{-\gamma}(X\circ(-\gamma))|_t = \nabla_{-\gamma'(a-t)}X = -\nabla_{\gamma'(a-t)}X = -D^\gamma_t(X\circ\gamma)|_t . $$
And so in general $D_t^{-\gamma}(X^-)=-D^\gamma_t(X)$ so clearly $X$ is parallel to $\gamma$ iff $X^-$ is parallel to $-\gamma$.

Then we have that $P^\gamma_{a,0}=P^{-\gamma}_{0,a}$ since if $X$ is the parallel transport of $v\in T_{\gamma(a)}M=T_{-\gamma(0)}M$ along $\gamma$ then $X^-$ is its parallel
transport along $-\gamma$ and $X(0)=X^-(a)$.
Thus we have
$$ (P^\gamma_{a,0})^{-1} = P^\gamma_{0,a} = P^{-\gamma}_{a,0} . $$
So parallel transport along $\Gamma(a-s,a)$ is the inverse of parallel transport along $\Gamma(s,a)$, and parallel transport along $\Gamma(0,a-t)$ is the inverse of parallel
transport along $\Gamma(0,t)$.

Now let $v\in T_pM$, and let $V$ be the vector field along $\Gamma$ obtained by parallel transport of $v$ first along $s\mapsto\Gamma(s,0)$ and then along $t\mapsto\Gamma(s,t)$.
Because parallel transport is the unique solution to a linear ODE and is thus smooth in its input, this is a smooth vector field along $\Gamma$.
And we have $V(0,0)=v$, since $s\mapsto V(s,0)$ is parallel along $\Gamma(s,0)$ we have that $D_sV(s,0)=0$.
And $t\mapsto V(s,t)$ is parallel along $t\mapsto\Gamma(s,t)$ so $D_tV(s,t)=0$ for all $s,t$.
We also know that
$$ R(x,y)v = (R(\partial_s\Gamma,\partial_t\Gamma)V)|_{(0,0)} = D_tD_sV(0,0) - D_sD_tV(0,0) = D_tD_sV(0,0) . $$
So we just need to compute $D_tD_sV(0,0)$.

Previously we showed that for a vector field $X$ along a curve $\gamma$,
$$ D_tX|_0 = \evdrv t_{t=0}(P^\gamma_{0,t})^{-1}(X(t)) = \lim_{\epsilon\to0}\frac{P^{\gamma(\epsilon-t)}_{0,\epsilon}(X(\epsilon)) - X(0)}\epsilon . $$
And so we have
$$ D_sV(0,t) = \lim_{\epsilon\to0}\frac{P^{\Gamma(\epsilon-s,t)}_{0,\epsilon}(V(\epsilon,t)) - V(0,t)}{\epsilon} , $$
and similarly
$$ D_tD_sV(0,0) = \lim_{\delta\to0}\frac{P^{\Gamma(0,\delta-t)}_{0,\delta}(D_sV(0,\delta)) - D_sV(0,0)}{\delta} . $$
Because linear maps are continuous, plugging in the expression for $D_sV(0,\delta)$ in the second expression gives
$$ D_tD_sV(0,0) = \lim_{\epsilon,\delta\to0}\frac{P^{\Gamma(0,\delta-t)}_{0,\delta}\circ P^{\Gamma(\epsilon-s,\delta)}_{0,\epsilon}(V(\epsilon,\delta)) - P^{\Gamma(0,\delta-t)}_{0,\delta}(V(0,\delta)) - P^{\Gamma(\epsilon-s,0)}_{0,\epsilon}(V(\epsilon,0)) + V(0,0)}{\delta\epsilon} . $$
Since by definition $s\mapsto V(s,t)$ and $t\mapsto V(t,0)$ are parallel we see that
$$ \eqalign{
    V(\epsilon,\delta)&=P^{\Gamma(\epsilon,t)}_{0,\delta}\circ P^{\Gamma(s,0)}_{0,\epsilon}v,\cr
    P^{\Gamma(0,\delta-t)}_{0,\delta}V(0,\delta)&=V(0,0)=v,\cr
    P^{\Gamma(\epsilon-s,0)}_{0,\epsilon}V(\epsilon,0)&=V(0,0)=v.
} $$
The first line is the definition, and the second two lines follow by $P^{-\gamma}$ and $P^\gamma$ being inverses.
Thus we see
$$ D_sD_tV(0,0) = \lim_{\epsilon,\delta\to0}\frac{P^{\Gamma(0,\delta-t)}_{0,\delta}\circ P^{\Gamma(\epsilon-s,\delta)}_{0,\epsilon}\circ P^{\Gamma(\epsilon,t)}_{0,\delta}\circ P^{\Gamma(s,0)}_{0,\epsilon}v- v}{\delta\epsilon} . $$
Taking this limit over the path $\epsilon=\delta$ we see that the parallel transport is just holonomy about $\gamma_\epsilon$,
$$ D_sD_tV(0,0) = \lim_{\epsilon\to0}\frac{\hol_{\gamma_\epsilon}v-v}{\epsilon^2} = \lim_{\epsilon\to0}\frac{\hol_{\gamma_{\sqrt\epsilon}}v-v}{\epsilon} = \evdrv t_{t=0}\hol_{\gamma_{\sqrt t}}(v) . $$

So putting this together we see that for all $v\in T_pM$,
$$ R(x,y)v = D_sD_tV(0,0) = \evdrv t_{t=0}\hol_{\gamma_{\sqrt t}}(v) $$
as desired.

\bprob

    Let $\pi\colon E\to N$ be a smooth vector bundle of rank $r$, and let $f\colon M\to N$ be a smooth map.
    Define
    $$ f^*E = M\times_NE $$
    to be the pullback, and let $\pi_f\colon f^*E\to M$ be given by projection onto $M$.
    The previous homework showed that $f^*E$ is a smooth manifold.
    \benum
        \item Given a local trivialization of $E$,
        $$ \Phi\colon \pi^{-1}U\to U\times\bR^r, $$
        construct a local trivialization of $f^*E$,
        $$ \tilde\Phi\colon \pi_f^{-1}(f^{-1}U) \to f^{-1}U \times\bR^r . $$

        \item Show that if $g_{\alpha\beta}$ are the transition functions of $E$, then the transition functions of $f^*E$ are $g_{\alpha\beta}\circ f$.
        Conclude that with these local trivializations and transition functions, $f^*E$ is a vector bundle of rank $r$ over $M$ with fibers
        $$ (f^*E)_p \cong E_{fp} . $$
        This is the {\emphcolor pullback bundle} of $E$.

        \item Prove that there is a natural bijection
        $$ \Gamma(f^*E) \cong \set{X\colon M \to E\hbox{ smooth}}[\pi\circ X=f] . $$
        In particular there is a natural bijection between sections of $f^*TN$ and vector fields along $f$.

        \item Let $f\colon M\to N$ be smooth.
        Show that there is a vector bundle
        $$ \hom(TM,f^*TN) \to M $$
        whose fibers are given by $\hom(T_pM,T_{fp}N)$.
        Then prove that $p\mapsto(p,df_p)$ defines a smooth section of $\hom(TM,f^*TN)$.
    \eenum

\eprob

\benum
    \item We have that
    $$ \pi_f^{-1}(f^{-1}U) = \set{(p,e)\in f^*E}[fp=\pi e\in U] . $$
    So we can define $\tilde\Phi\colon\pi_f^{-1}(f^{-1}U)\to f^{-1}U\times\bR^r$ by
    $$ \tilde\Phi(p,e) = (p,\proj_2\circ\Phi(e)) . $$
    We have that $\proj_1\circ\tilde\Phi=\pi_f$.
    And the fiber along $p$ of $f^*E$ is
    $$ (f^*E)_p = \set{(p,e)\in f^*E}[\pi e=fp] , $$
    which is isomorphic to $E_{fp}$ via $(p,e)\mapsto e$.
    For $p\in f^{-1}U$, restricting $\tilde\Phi$ to $f^*E_p$ gives us the map $(p,e)\mapsto(p,\proj_2\circ\Phi(e))$, and since $\proj_2\circ\Phi$ is a linear isomorphism of $E_{fp}$,
    so is $\proj_2\circ\tilde\Phi$.

    So we need only show that $\tilde\Phi$ is a diffeomorphism.
    We will explicitly construct an inverse, so a map $\tilde\Psi\colon f^{-1}U\times\bR^r\to\pi_f^{-1}(f^{-1}U)\subseteq f^*E$.
    To define a map into $f^*E$ it is sufficient to define maps into $E$ and $M$ which commute with $\pi$ and $f$.
    We give these by
    $$ \eqalign{
        \psi_1\colon f^{-1}U\times\bR^r\to M,\quad&(p,v)\mapsto p,\cr
        \psi_2\colon f^{-1}U\times\bR^r\to E,\quad&(p,v)\mapsto\Phi^{-1}(fp,v).
    } $$
    Note that $\psi_2$ is well-defined and smooth as the composition $(p,v)\mapsto(fp,v)\in U\times\bR^r\xmapsto{\Phi^{-1}}\psi_2(p,v)\in E$.
    And we have
    $$ \pi\circ\psi_2(p,v) = \pi\circ\Phi^{-1}(fp,v) = \proj_1(fp,v) = fp = f\circ\psi_1(p,v) , $$
    so by the universal property we have a smooth map $\tilde\Psi\colon f^{-1}U\times\bR^r\to f^*E$ by
    $$ \tilde\Psi(p,v) = (\psi_1(p,v),\psi_2(p,v)) = (p,\Phi^{-1}(fp,v)) . $$
    Notice that $\tilde\Psi$ maps into $\pi_f^{-1}(f^{-1}U)$ since
    $$ f\circ\pi_f\circ\tilde\Psi(p,v) = fp \in U , $$
    and as an open submanifold we can restrict $\tilde\Psi$'s codomain to it.

    Now we have that
    $$ \tilde\Psi\circ\tilde\Phi(p,e) = \tilde\Psi(p,\proj_2\circ\Phi(e)) = (p,\Phi^{-1}(fp,\proj_2\circ\Phi(e))) . $$
    And we have that since $\proj_1\circ\Phi=\pi$,
    $$ \Phi(e) = (\pi(e),\proj_2\circ\Phi(e)) = (fp,\proj_2\circ\Phi(e)) , $$
    we get
    $$ \tilde\Psi\circ\tilde\Phi(p,e) = (p,\Phi^{-1}\circ\Phi(e)) = (p,e) . $$
    And conversely
    $$ \tilde\Phi\circ\tilde\Psi(p,v) = \tilde\Phi(p,\Phi^{-1}(fp,v)) = (p,\proj_2\circ\Phi\circ\Phi^{-1}(fp,v)) = (p,v) $$
    as desired.

    So $\tilde\Psi=\tilde\Phi^{-1}$, showing that $\tilde\Phi$ is indeed a diffeomorphism and thus a local trivialization.

    \item Let $\Psi$ be another local trivialization of $E$ which gives us a local trivialization $\tilde\Psi$ of $f^*E$.
    Suppose that $\tau$ is the transition function between $\Phi,\Psi$ so
    $$ \Phi\circ\Psi^{-1}(q,v) = (q,\tau(q)v),\qquad q\in U,v\in\bR^r . $$
    Then we get for $p\in f^{-1}U,v\in\bR^r$,
    $$ \tilde\Phi\circ\tilde\Psi^{-1}(p,v) = \tilde\Phi(p,\Psi^{-1}(fp,v)) = (p,\proj_2\circ\Phi\circ\Psi^{-1}(fp,v)) = (p,\proj_2(q,\tau(fp)v)) = (p,\tau(fp)v) . $$
    So we see that the transition function between $\tilde\Phi$ and $\tilde\Psi$ is $\tilde\tau=\tau\circ f$ as desired.
    Thus $f^*E$ is indeed a vector bundle, and we showed before that its fibers are
    $$ (f^*E)_p \cong E_{fp} . $$

    \item By the universal property, a map $\sigma\colon M\to f^*E$ is the same as maps $\sigma_1\colon M\to E$ and $\sigma_2\colon M\to M$ such that $f\circ\sigma_2=\pi\circ\sigma_1$,
    where $\sigma_1=\pi_f^E\circ\sigma$ and $\sigma_2=\pi_f\circ\sigma$ ($\pi_f^E\colon f^*E\to E$ is projection onto $E$).
    In particular a section $\sigma\colon M\to f^*E$, so $\pi_f\circ\sigma=\id$, corresponds to just giving $\sigma_1=\pi_f^E\circ\sigma$.

    In particular if we have a vector field along $f$, so $X\colon M\to E$ such that $\pi\circ X=f$, then by the universal property this gives rise to a unique map $\hat X\colon M\to f^*E$ such
    that $X=\pi_f^E\circ\hat X$ and $\pi_f\circ\hat X=\id$.
    So $\hat X\in\Gamma(f^*E)$.
    This is summarized in the following diagram:

    \bigskip
    \centerline{\drawdiagram{
        $M$\cr
        &$F^*E$&$E$\cr
        &$M$&$N$
    }{
        \diagarrow{from={1,1}, to={2,3}, curve=1cm, text=$X$, y distance=.75cm}
        \diagarrow{from={1,1}, to={3,2}, curve=-1cm, text=$\id$, x distance=-.75cm}
        \diagarrow{from={1,1}, to={2,2}, dashed, text=$\hat X$, y distance=.25cm}
        \diagarrow{from={2,2}, to={2,3}, text=$\pi_f^E$, y distance=.25cm}
        \diagarrow{from={3,2}, to={3,3}, text=$f$, y distance=-.25cm}
        \diagarrow{from={2,2}, to={3,2}, text=$\pi_f$, x distance=-.25cm}
        \diagarrow{from={2,3}, to={3,3}, text=$\pi$, x distance=.25cm}
    }}

    Conversely if we have a section $\sigma\colon M\to f^*E$, then let $\bar\sigma=\pi_f^E\circ\sigma\colon M\to E$.
    This is a vector field along $f$ since
    $$ \pi\circ\bar\sigma = \pi\circ\pi_f^E\circ\sigma = f\circ\pi_f\circ\sigma = f , $$
    where the last equality is because $\sigma$ is a section.

    These maps $\hat\bul$ and $\bar\bul$ are inverses, since for $X\colon M\to E$ a vector field along $f$,
    $$ \bar{\hat X} = \pi_f^E\circ\hat X = X , $$
    where the last equality was explained above, and is due to the universal property.
    And if $\sigma\colon M\to f^*E$ is a section, then
    $$ \hat{\bar\sigma} = \hat{\pi_f^E\circ\sigma} . $$
    Since for any $X$, $\hat X$ is the unique map such that $X=\pi_f^E\circ\hat X$ and $\pi_f\circ\hat X=\id$, and since $\sigma$ satisfies these, we see that $\hat{\bar\sigma}=\sigma$
    as desired.

    Thus we have natural bijections
    $$ \Gamma(f^*E) \xtof{\quad\bar\bul\quad}[\quad\hat\bul\quad] \set{X\colon M\to E\hbox{ smooth}}[\pi\circ X=f] , $$
    as desired.
    In particular a local frame of $f^*E$ corresponds to $X_1,\dots,X_n\colon M\to E$ smooth vector fields along $f$ such that $X_i|_p\in E_{fp}$ are linearly independent, by following the isomorphism
    $f^*E_p\cong E_{fp}$ and $X_i\oto\hat X_i$.

    \item As stated in the recitation, given vector bundles on $M$ $E\to M,F\to M$, we can construct dual vector bundles $E^*\to M$ and tensor product bundles $E\otimes F\to M$
    whose fibers are $(E^*)_p=E^*_p$ and $(E\otimes F)_p=E_p\otimes F_p$ respectively.
    Putting this together we can construct a vector bundle $E^*\otimes F\to M$ whose fibers are $E^*_p\otimes F_p$, which is naturally isomorphic to $\hom(E_p,F_p)$.
    This gives a construction of $\hom(E,F)$ via the natural bijection
    $$ \coprod_{p\in M}\hom(E_p,F_p) \cong \coprod_{p\in M}E_p^*\otimes F_p = E^*\otimes F . $$
    In particular for $E=TM$ and $F=f^*TN$ the fibers are $\hom(E_p,F_p)=\hom(T_pM,T_{fp}N)$.
 
    We will construct the dual bundle and show that if $\sigma_1,\dots,\sigma_n$ is a local frame for $E\to M$, then $\epsilon^1,\dots,\epsilon^n$ is a local frame for $E^*\to M$
    where $\epsilon^1|_p,\dots,\epsilon^n|_p$ is the dual basis for $\sigma_1|_p,\dots,\sigma_n|_p$ relative to $E_p$.
    Similarly if $(\sigma_i),(\tau_j)$ are local frames for $E\to M,F\to M$ respectively on a common open subset, then $(\sigma_i\otimes\tau_j)$ is a local frame for $E\otimes F\to M$
    where $(\sigma_i\otimes\tau_j)|_p=\sigma_i|_p\otimes\tau_j|_p$ in $E_p\otimes F_p$.

    Then consider the map $p\mapsto(p,df_p)$.
    Since $f$ is smooth consider coordinates $(U,(x^i))$ of $M$ and coordinates $(V,(y^j))$ of $N$ such that $f(U)\subseteq V$.
    Then $\ipdv{y^j}\circ f$ is a vector field $M\to TN$ along $f$ (as the composition of the vector field $\ipdv{y^j}$ with $f$ it is smooth).
    This forms a local frame for $f^*TN$ considering dimensions.
    And let $dx^i$ denote the dual frame for $\ipdv{x^i}$ in $U$.
    Then $dx^i\otimes\ievpdv{y^j}_f$ forms a local frame for $T^*M\otimes f^*TN=\hom(TM,f^*TN)$.

    Since $df_p$ is linear $T_pM\to T_{fp}N$, we can represent it relative to this local frame, so
    $$ df = f^j_i\cdot dx^i\otimes\evpdv{y^j}_f . $$
    Recall that the identitification $\hom(V,W)\cong V^*\otimes W$ is given by mapping $T\colon V\to W$ to $a^j_i\phi^i\otimes w_j$ (where $\phi^i$ is the dual basis for $(v_i)\in V$ and $(w_j)$ is a basis for $W$)
    where $a^j_iw_j=T(v_i)$.
    Thus for $p\in U$,
    $$ f^j_i(p)\evpdv{y^j}_{fp} = df_p\parens{\evpdv{x^i}_p} = \pdvof{f^j}{x^i}(p)\cdot\evpdv{y^j}_{fp} , $$
    thus $f^j_i(p)=\pdvof{f^j}{x^i}(p)$.

    This means that $df$ has the local representation
    $$ df = \pdvof{f^j}{x^i}\cdot dx^i\otimes\evpdv{y^j}_f , $$
    and since its component functions $\ipdvof{f^j}{x^i}$ are smooth, this means that $df$ is a smooth section as desired.

    Now we actually construct $E^*$.
    Let $\pi_*\colon E^*\to M$ be the projection.
    Now let $\sigma_1,\dots,\sigma_n$ be a local frame for $E$, then let $\epsilon^1,\dots,\epsilon^n\colon M\to E^*$ be defined by $\epsilon^i|_p$ forming the dual basis for $\sigma_i|_p$.
    Now if we have another local frame $\tau_1,\dots,\tau_n$ for $E$ and $\delta^1,\dots,\delta^n$ be the dual local frame for $E^*$, we want to show $\epsilon^i,\delta^j$ are compatible.
    Let us write $\delta^i$ in terms of $\epsilon^j$ and $\sigma_i$ in terms of $\tau_j$,
    $$ \delta^i = a^i_j\epsilon^j,\qquad \sigma_i = b^j_i\tau_j . $$
    Evaluating at $p$ and composing with $\sigma_j|_p$ on the left and $\delta^j|_p$ on the right shows
    $$ a^i_j(p) = \delta^i|_p(\sigma_j|_p),\qquad b^j_i(p) = \delta^j|_p(\sigma_j|_p) . $$
    Thus $a^i_j=b^j_i$, and since $\sigma_i,\tau_j$ are compatible this means that $b^j_i$ is invertible and smooth and thus so is $a^i_j$, so $\epsilon^i,\tau^j$ are compatible.
    Since local frames for $E$ cover $M$, local dual frames for $E^*$ cover $M$, so $E^*$ is a smooth vector bundle of the same rank as $E$, and local frames of $E$ give rise to dual local frames of $E^*$.

    Similarly we construct $E\otimes F$.
    Let $\pi_\otimes\colon E\otimes F\to M$ be the projection.
    Similar to before we take local frames $\sigma_i,\alpha_j$ for $E,F$ defined on a common open subset, and construct local frames $\sigma_i\otimes\alpha_j$ for $E\otimes F$ whose value at $p\in M$ is $(\sigma_i\otimes\alpha_j)|_p=\sigma_i|_p\otimes\alpha_j|_p$.
    These are linearly independent at each point because the tensor of two bases form a basis for the tensor product.
    If we have other local frames $\tau_i,\beta_j$ for $E,F$, then we show that $\tau_i\otimes\beta_j$ and $\sigma_i\otimes\alpha_j$ are compatible.
    Suppose
    $$ \tau_i = a^k_i\sigma_k,\qquad \beta_j = b^\ell_j\alpha_\ell $$
    so that $a^k_i,b^\ell_j$ are the transition maps for the frames of $E$ and $F$ respectively, and are thus smooth.
    Then by evaluating at each point we see
    $$ \tau_i\otimes\beta_j = a^k_ib^\ell_j \sigma_k\otimes\alpha_\ell . $$
    Thus the transition functions between $\tau_i\otimes\beta_j$ and $\sigma_k\otimes\alpha_\ell$ are given by $a^k_ib^\ell_j$.
    As the product of smooth real-valued maps, each of these is smooth, and thus the frames are compatible.
    Thus we have shown that $E\otimes F$ is a vector bundle, and the tensor of local frames of $E$ and $F$ gives a local frame of $E\otimes F$ as desired.
\eenum



